\section{Route persistence, species multisections, and conjugation of Möbius type}
\label{sec:persistencia-ruta-mobius-familias}

This section specifies a distinction that, if formulated too
briefly, reverses the logic of the construction. The route is not a datum that
a particle chooses once its mass is known. Nor is it an auxiliary drawing
without causal content. It is the observable and oriented history of a surviving extension.
The extension additionally preserves compatible depth, the
register, carries, and memory that do not fit into a finite cutoff. In this
precise sense,
\[
 \boxed{\text{an HMT particle is the compatible persistence of an
 enriched route}.}
\]
The phrase and the definition by extensions express the same object at two
levels: the route is its visible history; the extension is the complete tower of
that history with all its compatible prolongations.

\subsection{Persistence as a section of an inverse system}

Let \((\mathcal S_m,\rho_{m+1,m})_{m\geq0}\) be the system of states that
survive to depth \(m\), and let
\[
 x=(x_m)_{m\geq0}\in\Omega_{\rm HMT}^{\rm surv}
 =\varprojlim_m\mathcal S_m.
\]
Route shadows satisfy the compatibility square
\[
 \operatorname{tr}_{m+1,m}\!\left(\operatorname{sh}_{m+1}(x_{m+1})\right)
 =\operatorname{sh}_{m}\!\left(\rho_{m+1,m}(x_{m+1})\right)
 =\operatorname{sh}_{m}(x_m).
\]
Therefore
\[
 \gamma_x=\bigl(\gamma_x^{[m]}\bigr)_{m\geq0},
 \qquad \gamma_x^{[m]}=\operatorname{sh}_m(x_m),
\]
is a compatible history, not a sequence retroactively fitted.

Over each cut \(\gamma_x^{[m]}\) one has the discrete bundle
\(\mathcal B_{\gamma_x^{[m]}}\). A persistent section is a family
\[
 s=(s_m)_{m\geq0},
 \qquad
 s_m\in\Gamma\!\left(\mathcal B_{\gamma_x^{[m]}}\right),
 \qquad
 \operatorname{res}_{m+1,m}(s_{m+1})=s_m.
\]
This equation is the mathematical content of persistence: each finite reading
extends without contradicting the preceding readings. The
particle tuple of \cref{chap:particula} adds occupation,
representation, and propagation, but does not alter this provenance.

\subsection{Internal species as an orbital quotient of the multisection}

In the nonadic atlas, each oriented cell \(c=(r,c')\in\mathcal C_{81}\)
carries a route family
\[
 \mathcal R(c)=
 (n_A,s_C,k_{\Delta_4},\nu_{120},\nu_{270},\tau_{12}).
\]
The image of \(\mathcal R\) contains fifty-six values. Once the
chart and its order are fixed, rank \(\kappa\in\{0,1,2,3,4\}\) within each fibre
produces the injective lift
\[
 c\longmapsto \bigl(\mathcal R(c),\kappa(c)\bigr).
\]
The five-state alphabet is minimal, since two fibers of
cardinality five exist. This selection is internal, prospective,
and blind to masses and particle names.

The complete census supporting this statement is
\[
 81=45_{\rm leptones}+36_{\rm quarks}
   =25_{\ell^-}+20_{\nu}+20_{d\text{-type}}+16_{u\text{-type}}.
\]
The fifty-six route fibers possess the histogram.
\[
 41\times1+10\times2+2\times3+1\times4+2\times5=81.
\]
Therefore, \(56\) counts route signatures, whereas \(81\) counts oriented
cells; memory \(\kappa\) distinguishes cells that one and the same signature does not
separate.

To classify internal species, this map is not inverted. The
quotient is taken.
\[
 q:\mathcal C_{81}\longrightarrow\Sigma_{13},
 \qquad
 q(c)=\bigl(\operatorname{familia}(c),G(c)\bigr),
\]
whose image comprises thirteen types. The output associated with a type \(y\) is the
finite multisection of the atlas
\[
 \boxed{
 \mathfrak S_y^{\rm atlas}=
 \{(\mathcal R(c),\kappa(c)):q(c)=y\}.}
\]
Their cardinals are
\[
 1,8,16;\qquad 2,4,6,8;\qquad 2,4,6,8;\qquad4,12,
\]
for charged leptons, neutrinos, and the two types of quark,
respectively. The species is the common type of that fiber; a particular particle
is a persistent section realized within it. Requiring the
name of a species by itself to select a cell destroys precisely the
degrees of orientation and conjugation that the fiber preserves.

The elevation to extensions is a different typed object. If
\(\pi_{81}:\Omega_{\rm HMT}^{\rm surv}\to\mathcal C_{81}\) denotes the cellular reader
of the declared enriched route, then
\[
 \widetilde{\mathfrak S}_y
 =\{x\in\Omega_{\rm HMT}^{\rm surv}:q(\pi_{81}x)=y\}
\]
is the corresponding enriched multisection. The finite certificate
proves the census and the obstruction in
\(\mathfrak S_y^{\rm atlas}\); its lift is relative to the reader
\(\pi_{81}\), not an automatic identification between a cell and a
deep history.

\begin{theorem}[Obstruction to the point selector and electronic exception]
Let \(\rho_{\rm c}(r,c')=(10-r,10-c')\) be central inversion. The thirteen fibers of
\(q\) are \(\rho_{\rm c}\)-stable. The level-zero charged-lepton fiber has
the unique fixed point, \((5,5)\). None of the other twelve fibers admits
a \(\rho_{\rm c}\)-equivariant point selector if \(\rho_{\rm c}\) acts trivially
on the coarse type.
\end{theorem}
\begin{proof}
The equation \(\rho_{\rm c}(r,c')=(r,c')\) is equivalent to \(r=c'=5\). The atlas census
places that cell in the fiber \((\ell^-,G0)\), which has cardinality one. Now let
\(y\neq(\ell^-,G0)\) and suppose that there exists an equivariant point section
\(a:\Sigma_{13}\to\mathcal C_{81}\). Since \(\rho_{\rm c} y=y\),
equivariance would require
\[
 a(y)=a(\rho_{\rm c} y)=\rho_{\rm c} a(y).
\]
Then \(a(y)\) would have to be a fixed point of \(\rho_{\rm c}\), but the unique fixed point belongs to the electron fiber. This is impossible. Therefore, in the other twelve fibers the natural output is a \(\rho_{\rm c}\)-orbit or a multisection; choosing a representative requires additional orientation or representation.
\end{proof}

The electron is exceptional in three compatible senses: \((5,5)\) is the
geometric unique center, the unique level-zero cell, and the center with respect
to which the raw lattice of differences is formed
\[
 \Lambda_{\rm raw}^{\Delta}
 =\bigl\langle A_{\rm raw}(c)-A_{\rm raw}(5,5):
 c\in\mathcal C_{81}\bigr\rangle_{\mathbb Z}.
\]
The raw signature of that cell is not null:
\[
 A_{\rm raw}(5,5)=(24,0,12,0,-12),
 \qquad \tau(5,5)=\texttt{+++---+++---}.
\]
The origin \(v_e=(0,0,0,0,0)\) of the lattice of action is a convention
relative to the electron and not the raw signature of the central cell.

The electronic algebraic invariant is evaluated by
\[
 \mathfrak e_0=
 \frac{\sqrt3}{4}
 \left(e^{\varphi/\pi^2}-22\alpha_{\rm HMT}^{3}\right)
 (1+15\Delta_4),
 \qquad
 \Delta_4=
 \frac{\pi-e\log\pi}{270}\left(1+\frac{\pi}{729}\right),
\]
\[
 \mathfrak e_0
 =0.510998922488066072509870877191349\ldots.
\]
The expression is dimensionless; the chart
\(\varsigma_{u_E}:y\mapsto y\,u_E\) subsequently introduces an energy unit.
The center \((5,5)\) and its uniqueness are exact finite results. The selection
of the coefficients \(22\) and \(15=3\cdot5\) is deduced, with explicit initial
structure, through the complement
of the five regions of \(\pi\) in \(K_e=\mathbb Z/9\mathbb Z\times\mathbb F_3\):
\(-22=-|K_e\setminus\iota(R_\pi)|\) and
\(+15=|R_\pi\times\mathbb F_3|\). The central raw signature retains a distinct type
and acts as an independent control of the basal selection.

The cell \((5,5)\) must not be confused with the energy signature
\(555555\): the latter belongs to the transversal representative of one of
the five regions of \(\pi\). They are distinct objects sharing the central
digit.

In the quark sector, the residual character
\[
 \operatorname{col}(r,c')=r-c'\pmod3
\]
refines the multisection into three classes of twelve cells and satisfies
\(\operatorname{col}(\rho_{\rm c} c)=-\operatorname{col}(c)\). This constructs the internal
color tritic. The physical representation of \(SU(3)\) is a
subsequent reader; it is not needed to prove the finite distribution.

\subsection{Stratification of extension, route, record, and signature}

The maps \(q\), \(\mathcal R\), the cellular projection, and the
realization reader distinguish four related and typed classifications:
\begin{enumerate}[label=(\roman*)]
 \item the thirteen internal family--level fibers of \(q\);
 \item the fifty-six route families of \(\mathcal R\);
 \item the eighty-one positions and their three hundred and twenty-four
       orientations;
 \item the individualized physical realizations by representation and codomain.
\end{enumerate}
The four layers are censused before the metrological contrast. For each class,
the comparison requires the tuple
\[
(\text{position},\text{fiber},\text{routes},\text{readers},
\text{record},\text{signature},\text{character},\text{towers},
\text{dimensional chart}).
\]
A reconstruction on a subset proves only that subcase and does not
replace any component of the tuple.
Given \(A\), \(C^*\), and the channels \(q_\pm\), the global tower belongs to the
infinite semigroup
\[
 s_n=q_-^n-q_+^n,
 \qquad
 n\in\langle90,120\rangle
 =\{90,120\}\cup\{30r:r\geq6\}.
\]
The heights \(90,120,180,210,240,270,360,420\) are finite representatives of the semigroup; the
preceding recurrence prolongs the complete tower. and the
dimensionless character of a signature is
\[
 \chi_{\rm mass}(v)=\exp\!\left(
 n_AA_{\rm rad}+\frac{s_C}{6}C^*_{\rm rad}
 +k\Delta_4+\nu_{120}s_{120}+\nu_{270}(135\Delta_4^2)
 \right).
\]
The moments are global: species weight them through their route, their
record, and their fibre operator. The readers \(K_D\) and \(K_\Omega\)
determine the bidirectional refinement before comparison, while the
integral table of species separately publishes the internal genealogy, the
dimensional chart, and the external comparison.

\subsection{Conjugation, anti-section, and Möbius band}

On a dodecaphasic word the exact involution is defined
\[
 C_M(\tau)_m=-\tau_{11-m},
 \qquad C_M^2=I.
\]
Linear functionals with symmetric weights are odd under \(C_M\),
whereas quadratic correlations compatible with reversion are
even. Thus, orientation and linear charge reside in the odd sector,
whereas quadratic memory resides in the even sector. This does not prove that
the complete exponential mass character is even: if
\(\chi(\tau)=e^{\ell(\tau)}\) and \(\ell\) is odd, then
\(\chi(C_M\tau)=\chi(\tau)^{-1}\).

The preceding involution acts exactly on words. To speak of an
anti-section, one additionally fixes a lift \(\widetilde C_M\) to the bundle of
sections that preserves admissibility, record, and truncations:
\[
 \operatorname{res}_{m+1,m}\widetilde C_M
 =\widetilde C_M\operatorname{res}_{m+1,m},
 \qquad
 \operatorname{sh}(\widetilde C_Mx)=C_M\operatorname{sh}(x).
\]
Only under that lift is the anti-section of \(s\)
\(\widetilde C_Ms\) on the inverted route.

The word ``Möbius'' has here a precise relative formulation. Let
\(\ell_{\rm ret}\) be a return loop of the enriched route that generates
\[
 \langle\ell_{\rm ret}\rangle\simeq\pi_1(S^1)\simeq\mathbb Z.
\]
This return group is not the finite phase \(C_9\). Suppose that its
lift to the
ledger has a sign character
\[
 \varepsilon_x:\langle\ell_{\rm ret}\rangle\longrightarrow\{\pm1\},
 \qquad \varepsilon_x(\ell_{\rm ret})=-1.
\]
The associated interval bundle is
\[
 \mathcal M_x=
 ([0,1]\times[-1,1])/(0,u)\sim(1,-u).
\]
Its first Stiefel--Whitney class is nontrivial; therefore
\(\mathcal M_x\) is a Möbius band. One circuit exchanges the two
local sheets, and two circuits restore the orientation. This does not require the
complete record to return to the same state: orientation return and memory
return remain distinct types.

The theorem is relative to the lift, the return loop, and the character
\(\varepsilon_x\), all of which are declared. It does not arise from observing two signs.
The reading of sheets additionally requires distinguishing the split grading \(R\)
from the conjugation \(C\) and imposing
\[
 CRC^{-1}=-R,
 \qquad CT(C^*)C^{-1}=T(-C^*).
\]
The particle--antiparticle reading ultimately requires the physical
representation to invert odd carriers, preserve even quadratic
carriers, and intertwine evolution and the complete character in the form declared
by the representation. Under those premises, particle and antiparticle are
the two conjugate orientations of a single persistence band.

\begin{theorem}[Even--odd decomposition on a double cover]
\label{pf84:A10-md-mobius}
Let \(p:\widetilde X\to X\) be a principal double cover with involution
\(\iota\). For every \(f\in C^0(\widetilde X,\mathbb R)\), the projectors
\[
 f^+=\frac12(f+\iota^*f),\qquad
 f^-=\frac12(f-\iota^*f)
\]
yield the decomposition \(f=f^++f^-\), with
\[
 f^+\circ\iota=f^+,\qquad f^-\circ\iota=-f^-.
\]
The even part descends to a function on \(X\); the odd part is a
section of the associated sign line. For the orientation cover
of the Möbius band, \(w_1\ne0\) obstructs a global continuous
nonvanishing section of that line.
\end{theorem}

\begin{proof}
The identities of the two projectors prove the sum and the parities. An invariant function is constant on the fibers of \(p\) and therefore descends; an anti-invariant function obeys the transition law of the sign representation. In the Möbius band, the associated line is nontrivial and its first Stiefel--Whitney class is nonzero, so every global continuous section has some zero. The conclusion excludes a global continuous section \emph{without zeros}; it does not deny the existence of sections with zeros.
\end{proof}

This classic lemma types the even--odd sector of the cover; by itself it identifies
neither word reversion, return holonomy, split
gradation, sheet conjugation, nor the parity of a
mass signature. Those identifications require their own intertwiners.

\subsection{Algebraic TRIT and temporal reader}

The declared quadratic realization associates an algebraic type with the visible sign of the TRIT. For
\(\bar\tau\in\{+1,0,-1\}\), let
\[
 \mathcal A_{\bar\tau}=
 \mathbb R[J]/(J^2+\bar\tau).
\]
Then
\[
 \begin{array}{c|c|c}
 \bar\tau&J^2&\text{regime}\ \\ \hline
 +1&-1&\text{elliptic/complex}\ \\
 0&0&\text{parabolic/nilpotent}\ \\
 -1&+1&\text{hyperbolic/split}.
 \end{array}
\]
The corresponding exponentials are
\[
 e^{tJ}=\cos t+J\sin t,
 \qquad e^{tJ}=1+tJ,
 \qquad e^{tJ}=\cosh t+J\sinh t,
\]
and in the split sector the projectors
\(P_\pm=(1\pm J)/2\) appear. These identities are verified exactly in the
declared realization. The
curvature does not arise from the isolated sign: it appears when APP supplies links and
two-cells and a connection is fixed. In the published mixed sum--product
polarization, over the \(81\) plaquettes, the oriented curvature satisfies
\[
 F(S,S)=F(P,P)=0,
 \qquad F(S,P)=+1,
 \qquad F(P,S)=-1\pmod9.
\]
The corresponding finite Stokes identity controls that the boundary of each
rectangle reproduces the sum of its plaquettes.

Ordering the transitions yields the tritic temporal reader:
\[
 \begin{aligned}
 +1&\longmapsto\text{return, memory and past},\\
  0&\longmapsto\text{torsional threshold and present},\\
 -1&\longmapsto\text{bidirectional opening and future}.
 \end{aligned}
\]
Its exact predecessor is the difference between phase return and state return:
\[
 g(k+9)=g(k),
 \qquad B_{k+9}\neq B_k\quad\text{in general}.
\]
In the corresponding geometric realization, one speaks of a spherical--volumetric sector, a flat threshold with a torsional ledger, and a hyperbolic--areal sector. This correspondence defines a Dimensional Mechanics reader subject to the map between TRIT and geometry; it neither redefines the algebraic TRIT nor by itself identifies a sign with a spacetime curvature. The expression ``time crystal'' names this architecture in HMT terminology; a dynamic realization of a time crystal additionally requires a Hamiltonian or a periodic excitation protocol, a robust subharmonic response, and stability.

\subsection{Exact support of the forward--return read}

Bidirectional exchange does not rest solely on a verbal image. The
integer angular coordinates
\[
 W_+=6n_A+s_C,
 \qquad W_-=6n_A-s_C
\]
are obtained by means of
\[
 \binom{W_+}{W_-}
 =\begin{pmatrix}6&1\\6&-1\end{pmatrix}
 \binom{n_A}{s_C},
 \qquad
 \operatorname{SNF}\!\begin{pmatrix}6&1\\6&-1\end{pmatrix}
 =\operatorname{diag}(1,12).
\]
The lattice therefore has index twelve. The involution
\(s_C\mapsto-s_C\) exchanges \(W_+\leftrightarrow W_-\) and, equivalently,
\(q_+\leftrightarrow q_-\). This is the exact algebraic substrate
of forward and return paths. Its interpretation as advanced and retarded propagation
requires the evolution operator of the physical representation.

The even projection used in the biangular reading satisfies, for every
function for which the expressions are defined,
\[
 \frac{f(\phi)+f(-\phi)}2=f_{\rm par}(\phi),
 \qquad
 \frac{(\pi-\phi)^2+(\pi+\phi)^2}{2}=\pi^2+\phi^2.
\]
They are exact symmetrization identities. Their reading in terms of a
theoretical physics model of the absorber belongs to the propagation functor, not to the
isolated algebraic identity.

\subsection{Inequivalent spinorial periodicities of \texorpdfstring{\(2\pi\)}{2 pi} and \texorpdfstring{\(4\pi\)}{4 pi}}

In the classical spinorial lift,
\[
 U(\theta)=\exp\!\left(-\frac{i\theta}{2}\sigma_z\right),
 \qquad U(2\pi)=-I,
 \qquad U(4\pi)=I.
\]
The equality is exact in the double cover \(SU(2)\to SO(3)\). Its application
to an HMT section requires specifying the spinorial representation of the
transport. Once fixed, it explains why one turn changes the sign of the
section and two turns restore it. If the CAR relations are adopted in the
occupation algebra, exclusion follows from them; neither CAR nor Pauli is
inferred from the sign of \(2\pi\) without that additional structure.

The HMT interpretation ``areal sheet \(2\pi\), volumetric sheet \(4\pi\)'' is another object. Section~\ref{sec:bisagra-hojas-pi-phi2} formalizes its core in topological terms: one circuit of the base closes the projected boundary in \(2\pi\) and ends on the conjugate sheet, whereas the squared loop closes the oriented cover in \(4\pi\). The theorem is relative to the declared lift, holonomy, and geometric reader. Its further identification with a spinorial representation still requires the physical functor and is not deduced from their sharing the same factors.

\subsection{Map between internal multisections and physical representations}

The particle object, the projection from extension to route, the internal selector of the 81 cells, the electron center, the thirteen finite family multisections, their lift relative to the cellular reader \(\pi_{81}\), the color trit, the register, the signature, the character, and the towers form the composition
\[
\text{extension}\longrightarrow\text{route}\longrightarrow
\text{multisection}\longrightarrow\text{signature}\longrightarrow
\text{character}\longrightarrow\text{towers}.
\]

The remaining obligation is to construct a realization functor that pairs certain individualized and
composite physical labels with the relevant enriched multisections and reproduces the
declared refined signatures without consulting target masses or signatures. For
composite states, the register is composed before being projected onto five
coordinates. The domain of the functor preserves the central electron, the atlas of
families, and the internal character law.


\paragraph{Evidentiary typing of persistence, atlas and families.}
Route persistence, anti-sections, the bidirectional reading, and the tritic temporal reader fix the starting architecture. The inverse system, the internal selector, and the separation of types realize that architecture through declared morphisms. The atlas censuses and the identities of \(C_M\) on words hold throughout the finite system; their lift to sections and the Möbius band are proved from explicit starting structure. The electron coefficients \((-22,+15)\) are deduced from the same explicit starting structure. If a lift first opens a metrological difference or a target signature, its value is a calibrated reconstruction.
